% !TEX root = ../memoria.tex
\chapter{Validación y Resultados}
\label{cap:capitulo5}

\section{Alcance de la evaluación}
\label{sec:escenario}

La evaluación distingue entre pruebas de software, comparación entre motores de
propagación y validación empírica. La comprobación central de integración compara
Java y MATLAB en perfiles sintéticos con el mismo algoritmo FSPL+Deygout y supuestos
alineados: las diferencias observadas son inferiores a $10^{-8}$~dB. Una comparación
histórica recalculó con Java puntos de escenarios MATLAB con hipótesis distintas; no
se pudo repetir en esta integración sin el GeoTIFF histórico. No se dispone de una campaña de medidas de
campo trazables para validar las predicciones frente a RSRP real; por tanto, los resultados
de esta sección no se presentan como validación empírica.

El conjunto histórico de referencia incluye tres JSON MATLAB: \texttt{URJC\_Fuenlabrada\_Suburbano}
(3,5~GHz), \texttt{Sierra\_Guadarrama\_Rural} (0,7~GHz) y
\texttt{Madrid\_Centro\_Urbano} (3,5~GHz). Los tres declaran el modelo
\texttt{Longley-Rice} y el terreno \texttt{SRTM\_90m\_global}. La prueba
\texttt{MatlabScenarioComparisonTest} utiliza estos ficheros y el GeoTIFF local
\texttt{srtm\_fuenlabrada.tif} para los recálculos Java de aquella ejecución. El raster
histórico examinado medía
$3601\times3601$ muestras, está georreferenciado en WGS~84 y declara una escala de un
segundo de arco por píxel (resolución nominal cercana a 30~m, dependiente de la latitud).
El fichero lleva una etiqueta SRTM, pero no se conserva una ficha de adquisición o una
suma de verificación que permita documentar su cadena de procedencia.

\paragraph{Tipo de pixel.}
La clave GeoKey \texttt{GTRasterTypeGeoKey} del GeoTIFF histórico indicaba
\texttt{PixelIsPoint} en el raster de $3601\times3601$ píxeles. La corrección del
tratamiento \texttt{PixelIsArea} se cubrió con un raster sintético de $2\times3$ píxeles
y una prueba de regresión \texttt{PixelIsPoint}; por ello, no altera el tile local ni las
métricas históricas recogidas en el Cuadro~\ref{tab:comparacion_motores}.

\section{Pruebas del software}
\label{sec:resultados_web}

El informe Maven de la fase~1, fechado el 21/11/2026, contabilizó 47 pruebas: 46 ejecutadas y
una comparación Java--MATLAB optativa omitida, sin fallos ni errores. El XML JaCoCo de esa fase
(\texttt{backend/target/site/jacoco/jacoco.xml}, generado a las 21:22) registró 409 de
453 líneas cubiertas (90,3~\%) y 137 de 192 ramas (71,4~\%) en 26 clases de producción.
La cobertura de código mide ejecución de instrucciones; no demuestra la exactitud
radioeléctrica del motor.

En esa fase anterior, la compilación de producción del frontend Angular terminó correctamente. Se observaron
avisos del presupuesto de tamaño del paquete inicial y de dependencias CommonJS. Una
prueba funcional local en navegador, con seis JSON de prueba aislados, verificó la creación
y visualización del mapa, el borrado confirmado y cancelado, el tratamiento de JSON
inválidos, las respuestas GET tardías o fallidas y limpiar una comparación A=B. No es una
suite de navegador integrada en CI. Tampoco se midieron latencia de simulación, tiempo de
carga del mapa o rendimiento con usuarios concurrentes; esos requisitos siguen pendientes
de medición.

En la validación del candidato integrado \texttt{f8ae98775c03094f593bd3dfa852a0f8577cba7c},
Maven registró 82 casos: 79 ejecutados y superados y tres omitidos, sin fallos.
Los casos omitidos fueron el de SRTM local, el de recomendación con raster local y
la comparación histórica; no se consideran pruebas ejecutadas. Un control selectivo
adicional de Java--MATLAB con perfiles sintéticos fue satisfactorio, con diferencias
inferiores a $10^{-8}$~dB. El frontend superó siete pruebas y su compilación fue
correcta. Las cifras de cobertura JaCoCo citadas arriba pertenecen únicamente a la
fase~1 y no se extrapolan al candidato. En una revisión posterior se amplió la
suite del frontend a nueve pruebas, incluidas las regresiones de importaciones
obsoletas de calibración.

\section{Comprobación sintética de FSPL y pérdidas de sistema}
\label{sec:comprobacion_fspl_sintetica}

Una revisión independiente ejecutó MATLAB 25.2.0.3150157 (R2025b Update~4) con dos
\textit{sites} cartesianos en $[0;0;0]$ y $[1000;0;0]$~m, antenas TX y RX explícitamente
isotrópicas, frecuencia de 3,5~GHz y \texttt{rx.SystemLoss=0}. La función \texttt{fspl}
devolvió 103,329144109~dB. Con
\texttt{propagationModel('freespace')}, \texttt{Map='none'} y \texttt{sigstrength}, se
obtuvieron los resultados del Cuadro~\ref{tab:matlab_balance_sintetico}.

\begin{table}[H]
\centering
\caption{Comprobación sintética del balance de potencia en MATLAB.}
\label{tab:matlab_balance_sintetico}
\small
\begin{tabularx}{\textwidth}{|X|X|r|}
\hline
\textbf{Configuración TX} & \textbf{Configuración RX} & \textbf{\texttt{sigstrength} (dBm)} \\
\hline
53~dBm, \texttt{SystemLoss=0} & Isotrópica, \texttt{SystemLoss=0} & $-50{,}329144109$ \\
55~dBm, \texttt{SystemLoss=2} & Isotrópica, \texttt{SystemLoss=0} & $-50{,}329144109$ \\
\hline
\end{tabularx}
\end{table}

La diferencia entre ambas salidas fue 0,000000000~dB. Este resultado confirma la
consistencia del balance y la aplicación de \texttt{SystemLoss} en esa configuración
sintética; no demuestra qué antena efectiva usaron los escenarios históricos ni corrige
sus JSON. \texttt{sigstrength} devuelve potencia recibida total simulada: el ejercicio no
calcula SS-RSRP y no contiene medidas de campo. Fue una única ejecución; no se modificaron
los generadores ni los datos.

Como contraste de implementación, el \texttt{PropagationEngine} Java actual se ejecutó
con perfil sintético plano de 0~m, TX y RX a 10~m, distancia horizontal de
1\,000{,}000000000~m, frecuencia de 3,5~GHz y balance TX de
$40+15-2=53$~dBm. Java obtuvo FSPL de 103,329144108889~dB, difracción Deygout de 0~dB y
un valor proxy de $-79{,}316395924784$~dBm. Como referencia, el resultado MATLAB de
\texttt{sigstrength} ($-50{,}329144109$~dBm de potencia total) normalizado con
$10\log_{10}(66\times12)=28{,}987251815895$~dB da
$-79{,}316395924895$~dBm; la diferencia Java--MATLAB normalizado es
$+1{,}11\times10^{-10}$~dB. El motor Java no expone potencia total recibida en el mismo
formato, por lo que su valor es un proxy, no una medición ni un cálculo normativo de
SS-RSRP según 3GPP TS~38.215~\cite{3gpp_ts38215_2022}.

Este control sintético de un único enlace de línea de vista muestra concordancia numérica
entre la implementación Java y la referencia MATLAB normalizada bajo estas hipótesis; la
diferencia residual queda por debajo de la resolución publicada de la salida MATLAB y no
debe interpretarse como precisión significativa. El caso no permite atribuir ni descartar
causas para los RMSE históricos. Las diferencias de configuración de antena y potencia,
normalización por banda, terreno, modelo Longley--Rice y filtrado de puntos son hipótesis
que habría que aislar, no causas probadas por este control. Tampoco se validan la antena
efectiva de los JSON históricos, un terreno real ni predicciones con medidas de campo.

\section{Comparación entre los motores Java y MATLAB}
\label{sec:comparativa}

El motor Java calcula FSPL más difracción de Deygout sobre el perfil de elevación del
GeoTIFF local. En la ejecución histórica, para cada punto guardado en un JSON MATLAB,
el test volvió a calcular Java
en las mismas coordenadas y acumula la diferencia
$\Delta=\mathrm{RSRP}_{Java}-\mathrm{RSRP}_{MATLAB}$. Se informa el sesgo medio con signo
y la raíz del error cuadrático medio (RMSE), en dB. Los puntos que el motor no puede
evaluar por quedar fuera de la cobertura del raster local se excluyen y se cuentan por
separado.

\begin{table}[H]
\centering
\caption{Comparación histórica entre el motor Java y los escenarios MATLAB almacenados.}
\label{tab:comparacion_motores}
\small
\begin{tabularx}{\textwidth}{|X|r|r|r|r|}
\hline
\textbf{Escenario de referencia} & \textbf{Pares} & \textbf{Excluidos} &
\textbf{Sesgo medio (dB)} & \textbf{RMSE (dB)} \\
\hline
\texttt{Sierra\_Guadarrama\_Rural} & 3168 & 1683 & 12,389 & 19,830 \\
\texttt{URJC\_Fuenlabrada\_Suburbano} & 3598 & 0 & 7,866 & 12,242 \\
\texttt{Madrid\_Centro\_Urbano} & 2451 & 0 & -3,938 & 16,463 \\
\hline
\end{tabularx}
\end{table}

Los valores corresponden a una ejecución anterior, no repetida en la integración
actual por falta del GeoTIFF histórico. Muestran diferencias notables entre las
salidas de ambos motores: el sesgo
cambia de signo entre escenarios y los RMSE observados se sitúan entre 12,242 y
19,830~dB. El mayor número de exclusiones se produce en el escenario de la Sierra. La
prueba histórica identificó puntos que el motor no podía evaluar por la cobertura del raster
local; con el GeoTIFF examinado, que no declara celdas NoData, se atribuyen a ubicaciones
fuera de su extensión. Sin embargo, la revisión independiente ha señalado que debe
reconciliarse el tratamiento de puntos excluidos antes y después de exportar los escenarios;
esa equivalencia sigue bajo revisión y no se da por resuelta aquí. La prueba histórica
exigía al menos un punto evaluable y registraba métricas, sin definir un umbral de
aceptación; el procedimiento integrado actual añade tratamiento de valores recortados,
por lo que sus cifras no deben identificarse con las de esta tabla sin repetirlo.

Esta comparación no aísla el efecto de los modelos. Los JSON MATLAB declaran
\texttt{SRTM\_90m\_global}, mientras que Java utiliza el GeoTIFF local de un segundo de
arco; además, Longley--Rice y FSPL+Deygout son algoritmos distintos. También hay diferencias
nominales en el balance de potencia: MATLAB pasa $P_{tx}$ al modelo sin sumar ganancia de
antena ni restar pérdidas de cable, mientras Java usa EIRP; el desplazamiento nominal es
13~dB en URJC y Madrid y 16~dB en Sierra. Para Sierra a 700~MHz, los factores heredados son 52$\times$12 en MATLAB
(aproximadamente 27,96~dB) y 66$\times$12 en Java (28,99~dB), con una diferencia cercana a
1,03~dB. No representan una configuración documentada de ancho de banda o SCS. En el código
actual, Java construye una rejilla inclusiva con
$\lceil\text{área}/\text{paso solicitado}\rceil+1$ puntos por eje y guarda los pasos
solicitado y efectivo; el generador MATLAB calcula el paso efectivo como
$\text{área}/(n_\text{puntos}-1)$ y lo incluye en sus metadatos. Sin embargo, los JSON
históricos usados en esta comparación no se regeneraron, por lo que no incorporan estas
correcciones del generador ni constituyen una nueva comparación con artefactos armonizados.
Siguen pendientes de reconciliación el terreno efectivo y el filtrado de puntos antes y
después de exportar los escenarios. No se consideran resueltos esos límites ni se atribuye
el RMSE exclusivamente al modelo. Los valores siguen siendo una comparación entre
simulaciones, no una validación con mediciones de campo ni una prueba de conformidad del
estándar.

\section{Limitaciones identificadas}
\label{sec:limitaciones}

\begin{itemize}
  \item \textbf{Modelos no equivalentes}: los escenarios históricos MATLAB usan
    Longley--Rice y Java combina FSPL con difracción de Deygout. La fase~1 no modela
    edificios, reflexiones urbanas, interiores, patrones de antena ni interferencia; las
    extensiones locales posteriores no suplen todos estos fenómenos.
  \item \textbf{Alcance de un nodo}: esta revisión se centra en una antena y en su
    cobertura estimada; la validación empírica sigue pendiente.
  \item \textbf{Terreno}: el raster local presenta tags WGS~84 y escala de un segundo
    de arco; no se dispone de metadatos de adquisición para certificar su procedencia.
    Los JSON MATLAB conservan su propia declaración
    \texttt{SRTM\_90m\_global}, que no se debe sustituir por la etiqueta Java.
  \item \textbf{Límites de cobertura}: el backend consulta la extensión del raster y rechaza
    perfiles fuera de ella o con elevación no válida; la interfaz muestra avisos basados en
    la extensión y el área solicitada. La comprobación de la interfaz es parcial y no se
    presenta como una gestión completa de límites geográficos.
  \item \textbf{Historial}: el borrado confirmado, la comprobación de identidad del JSON y
    la reconciliación de selecciones están cubiertos por pruebas automatizadas y una prueba
    funcional local en navegador. No se dispone de una suite de navegador en CI. La
    persistencia se basa en ficheros, sin base de datos ni historial versionado
    completo. Los resultados de calibración vinculan las observaciones compatibles
    con el identificador de simulación y el hash de la predicción, y conservan
    las observaciones originales; esto no acredita una campaña de medidas reales.
  \item \textbf{Cartografía}: el catálogo activo ofrece tres fuentes raster externas. La
    integración PMTiles no tiene un tema activo y no se ha verificado el uso sin conexión.
  \item \textbf{Validación y rendimiento}: no se han incorporado medidas de campo; tampoco
    se han medido tiempos de simulación, carga del mapa, escalabilidad o usabilidad con
    usuarios. El test de comparación no fija un criterio de error aceptable.
\end{itemize}

En consecuencia, el prototipo alcanza los niveles descriptivo y predictivo interactivo de
la burbuja virtual, pero todavía no está sincronizado ni calibrado con una burbuja física.
El avance verificable consiste en generar, almacenar, consultar y representar predicciones
con dos motores diferenciados, con límites y precisión que deben interpretarse según las
evidencias expuestas.

\section{Avances locales posteriores a la fase 1}
\label{sec:avances_posteriores}

Esta sección recoge pruebas locales posteriores a la fase~1, incluida una integración
anterior identificada por el árbol Git \texttt{71df5ad5c1cce8d9c65863b1ce46a0cac121f018},
el árbol reorganizado \texttt{08d3a7e6fff368e7cd392e465c272c67ca88660d} y el
candidato integrado validado en la Sección~\ref{sec:resultados_web}.
Las cifras de pruebas y cobertura de la sección anterior corresponden a la fase~1,
salvo donde se identifica expresamente el candidato integrado.

Los ejemplos numéricos de las capacidades de apoyo en esta sección proceden de una
integración local anterior y no se atribuyen a una nueva ejecución del candidato.
Estas comprobaciones no establecen precisión frente a medidas de campo.

\paragraph{Recorrido integrado en navegador.}
En la integración reorganizada se comprobó el recorrido real del navegador para
colocar la antena, solicitar una simulación persistida y visualizar la cobertura,
utilizando un GeoTIFF y datos sintéticos aislados. No se utilizaron medidas de
campo ni se acredita con ello precisión del modelo en un entorno real.

\paragraph{Comparación controlada Java--MATLAB.}
Además del contraste histórico de la Sección~\ref{sec:comparativa}, se ensayaron perfiles
sintéticos plano y con cresta con el mismo algoritmo FSPL+Deygout y supuestos alineados
en ambas implementaciones. El control Java utilizó un JSON sintético generado previamente
por MATLAB; no se ejecutó MATLAB de nuevo para aquella comprobación. En la integración
reorganizada se comprobó también el contraste sintético con diferencia inferior a
$10^{-8}$~dB. Los proxies anteriores obtenidos fueron aproximadamente
$-79{,}316395925$ y $-111{,}418102426$~dBm, respectivamente, con diferencias inferiores
a $10^{-8}$~dB entre implementaciones. El control comprueba concordancia numérica para
esos perfiles, no independencia de modelos, exactitud en terreno real ni validación
frente a campo; no elimina las discrepancias de los escenarios históricos.

\paragraph{Calibración sintética.}
Se probó un ajuste con datos sintéticos, separando el conjunto usado para ajustar del
conjunto de reserva utilizado para evaluar. El emparejamiento utiliza el punto de rejilla
más cercano dentro del dominio y a una distancia máxima de $0{,}75$ veces el espaciado,
limitada a 100~m; se descartan puntos de reserva compartidos con el ajuste. El
desplazamiento es la media de observación menos predicción en el conjunto de ajuste.
En un ejemplo sintético, el desplazamiento obtenido fue $+3$~dB y, en la reserva, el
sesgo pasó de $-3$ a $0$~dB y el RMSE de 3,109126 a 0,816497~dB. Es una comprobación
del procedimiento, no una demostración de mejora sobre mediciones reales ni de
generalización. La calibración ajusta el \textit{proxy} de potencia guardado, no SS-RSRP
físicamente equivalente.

\paragraph{Recomendación de emplazamientos.}
Un prototipo local prueba una selección voraz (\textit{greedy}) explicable sobre rejillas
acotadas para un área cuadrada o un polígono simple dibujado, incluidos polígonos
cóncavos. En el caso poligonal se calculan rejillas sobre su envolvente y se conservan
los candidatos y objetivos interiores o en el contorno. Cada paso elige el candidato
que aporta más objetivos todavía no cubiertos. Se rechazan geometrías inválidas y se
limita el coste antes de calcular. La cobertura se refiere a puntos de la rejilla, no a
superficie continua. La recomendación usa antenas isotrópicas y no incorpora el patrón
direccional ni el \textit{clutter} manual.

En un ensayo simulado cuadrado con una tesela local se cubrieron 64 de 64 objetivos
con un nodo. En otro ensayo local, el dibujo y cálculo de un polígono en navegador
produjo 10 candidatos, 36 objetivos y 360 enlaces evaluados; un nodo cubrió los 36
objetivos. Estos casos no prueban un óptimo global ni rendimiento general en otros
escenarios. La evaluación completa de sus interacciones sigue pendiente.

\paragraph{Complementos radioeléctricos.}
Se probaron un patrón direccional opcional por transmisor, una pérdida de
\textit{clutter} introducida manualmente y la representación de un perfil con primera
zona de Fresnel.
Las pruebas radioeléctricas emplearon terreno sintético; el \textit{clutter} manual no
es un mapa de usos del suelo y no se verificó el perfil con un raster real. La
integración local posterior no acredita por sí misma la precisión de las predicciones.

\paragraph{Empaquetado local.}
Una prueba anterior con un GeoTIFF sintético devolvió indebidamente
\texttt{terrain\_data=SRTM\_30m\_local}. En el código integrado la etiqueta por
defecto es genérica y Compose configura explícitamente la procedencia sintética.
Se comprobó la configuración de Docker Compose, pero no se pudieron construir ni
ejecutar los contenedores por falta de espacio: la demo contenerizada y su
verificación visual siguen pendientes.

\paragraph{Comprobaciones pendientes.}
Quedan por confirmar la cobertura de pruebas del candidato y la reproducción de los
experimentos anteriores sobre el conjunto definitivo. Faltan la demo
Docker contenerizada, la evaluación completa de interacciones, la comparación histórica
con el raster original y una campaña de medidas de campo trazables con criterios de
aceptación. No se asigna a estas tareas un
resultado hasta ejecutarlas.
